% Katalog D: Tabellen und Gitter (Palette Slot 4).
\ZSFNewColumn
\StartChapter[kat:tabellen]{D · Tabellen}

\begin{runintext}
  Eine Tabellen-Umgebung mit fünf Reglern, dazu das kompakte
  \ZSFkeyword{valuegrid}. Alle Einträge tragen denselben Zellinhalt, damit
  die Regler und nicht die Daten den Unterschied machen.
\end{runintext}

\SubsectionBar[kat:tab-basis]{Kopfzeile und Zebra}
\ZSFindex{ZSFtable}
\ZSFindex{header}
\ZSFindex{zebra}

\begin{tablebox}[D-01 · ZSFtable]
  \begin{ZSFtable}{Y{1.0} Y{1.0}}
    \ZSFheaderRow \ZSFhead{Erste} & \ZSFhead{Zweite} \\
    Mustertext & zum Formvergleich \\
    Zweite Zeile & zeigt das Zebra \\
    Dritte Zeile & zeigt den Wechsel
  \end{ZSFtable}
\end{tablebox}
\Zweck{Vorbelegung: Kopfzeile an, Zebra ab Zeile 2. Kopfzeile braucht immer
beide Makros — ZSFheaderRow färbt, ZSFhead setzt die Schrift.}

\begin{tablebox}[D-02 · header=false]
  \begin{ZSFtable}[header=false]{Y{1.0} Y{1.0}}
    Mustertext & zum Formvergleich \\
    Zweite Zeile & Zebra ab Zeile 1 \\
    Dritte Zeile & zeigt den Wechsel
  \end{ZSFtable}
\end{tablebox}
\Zweck{Ohne Kopfzeile beginnt das Zebra bei der ersten Zeile — für reine
Datenlisten.}

\begin{tablebox}[D-03 · zebra=false]
  \begin{ZSFtable}[header=false, zebra=false]{Y{1.0} Y{1.0}}
    Mustertext & zum Formvergleich \\
    Zweite Zeile & ohne Streifen \\
    Dritte Zeile & zeigt den Wechsel
  \end{ZSFtable}
\end{tablebox}
\Zweck{Keine Streifen — für sehr kurze Tabellen.}

\SubsectionBar[kat:tab-mass]{Schrift, Zeilenhöhe, Polsterung}
\ZSFindex{rows}
\ZSFindex{colsep}

\begin{tablebox}[D-04 · font=dense]
  \begin{ZSFtable}[font=dense]{Y{1.0} Y{1.0}}
    \ZSFheaderRow \ZSFhead{Erste} & \ZSFhead{Zweite} \\
    Mustertext & zum Formvergleich \\
    Zweite Zeile & eine Stufe kleiner
  \end{ZSFtable}
\end{tablebox}
\Zweck{Für lange Register mit knappem Zellinhalt. Derselbe Reglername wie
an der Box, dieselbe Bedeutung.}

\begin{tablebox}[D-05 · rows=roomy]
  \begin{ZSFtable}[rows=roomy]{Y{1.0} Z{1.0}}
    \ZSFheaderRow \ZSFhead{Erste} & \ZSFhead{Zweite} \\
    Mustertext & $\dfrac{\Delta y}{\Delta x}$ \\
    Zweite Zeile & $\dfrac{a+b}{2}$
  \end{ZSFtable}
\end{tablebox}
\Zweck{Mehr Luft, wo hoher Zellinhalt sonst an die Zeilenkante stösst.
Vorbelegung: normal.}

\begin{tablebox}[D-06 · rows=tight]
  \begin{ZSFtable}[rows=tight, header=false]{Y{1.0} Q{1.0}}
    Mustertext & zum Formvergleich \\
    Zweite Zeile & knapp \\
    Dritte Zeile & dicht
  \end{ZSFtable}
\end{tablebox}
\Zweck{Knappe Zeilen für Register mit einzeiligem Kurzinhalt.}

\begin{tablebox}[D-07 · colsep=tight]
  \begin{ZSFtable}[font=dense, colsep=tight]{Y{1.4} Z{0.9} Z{0.9} Z{0.8}}
    \ZSFheaderRow \ZSFhead{Erste} & \ZSFhead{A} & \ZSFhead{B} & \ZSFhead{C} \\
    Mustertext & 12 & 24 & 36 \\
    Zweite Zeile & 20 & 40 & 60
  \end{ZSFtable}
\end{tablebox}
\Zweck{Knappe Zellpolsterung für vielspaltige Tabellen — zugleich zwischen
den Spalten und an der Boxkante. Vorbelegung: normal.}

\SubsectionBar[kat:tab-struktur]{Zellverbund, Bild, Linien}
\ZSFindex{ZSFspan}
\ZSFindex{ZSFimage}

\begin{tablebox}[D-08 · ZSFspan]
  \begin{ZSFtable}{Y{1.2} Z{0.9} Z{0.9} Z{0.8}}
    \ZSFheaderRow \ZSFhead{Erste} & \ZSFspan{C}{2}{\ZSFhead{Gruppe}}
                                  & \ZSFhead{C} \\
    Mustertext & 12 & 24 & 36 \\
    \ZSFspan{L}{4}{\ZSFlabel{Abschnittszeile über die volle Breite}} \\
    Zweite Zeile & 20 & 40 & 60
  \end{ZSFtable}
\end{tablebox}
\Zweck{Ein Baustein für drei Rollen: gruppierter Kopf, Formel über die
Breite, Abschnittszeile. Die Ausrichtung ist Pflichtargument, der Inhalt
bricht nicht um.}

\begin{tablebox}[D-09 · ZSFimage in der Zelle]
  \begin{ZSFtable}[colsep=tight]{Z{0.8} Y{1.0} Y{1.2}}
    \ZSFheaderRow \ZSFhead{Bild} & \ZSFhead{Erste} & \ZSFhead{Zweite} \\
    \ZSFimage{example-image-a} & Mustertext & ohne Höhenangabe \\
    \ZSFimage{example-image-b} & Zweite Zeile & gleich hoch von selbst
  \end{ZSFtable}
\end{tablebox}
\Zweck{Das Bild ohne eigene Box; die Höhe stellt der Container, nicht der
Aufruf.}

\begin{tablebox}[D-10 · Linien]
  \begin{ZSFtable}[header=false, zebra=false]{Y{1.0} | Q{1.0}}
    Mustertext & zum Formvergleich \\ \hline
    Zweite Zeile & nach hline \\
    Dritte Zeile & senkrecht getrennt
  \end{ZSFtable}
\end{tablebox}
\Zweck{Vorbelegung ist ohne Linien, aber hline und der Pipe in der
Spaltenangabe bleiben erlaubt — nur ihre Farbe wird nicht lokal gesetzt.}

\SubsectionBar[kat:tab-spalten]{Spaltentypen}
\ZSFindex{Spaltentyp}

\begin{tablebox}[D-11 · L C R F]
  \begin{ZSFtable}[header=false, zebra=false]{L C R F{1}}
    L links & C zentriert & R rechts & F{1}: x^2+1
  \end{ZSFtable}
\end{tablebox}
\Zweck{Gleichverteilt mit Faktor 1; F ist der neutrale Typ für Mathe-Inhalt.}

\begin{tablebox}[D-12 · Y Z Q mit Gewicht]
  \begin{ZSFtable}[header=false, zebra=false]{Y{1.4} Z{0.8} Q{0.8}}
    Y links, Gewicht 1.4 & Z zentriert & Q rechts
  \end{ZSFtable}
\end{tablebox}
\Zweck{Dieselben drei Ausrichtungen mit frei wählbarem Gewicht; die Summe
der Faktoren entspricht der Anzahl Spalten.}

\SubsectionBar[kat:tab-grid]{Wertegitter}
\ZSFindex{valuegrid}

\begin{valuegrid}[D-13 · valuegrid]{3}
  Mustertext & 12 & 24 & 36 \\
  Zweite Zeile & 20 & 40 & 60 \\
  Dritte Zeile & 30 & 60 & 90
\end{valuegrid}
\Zweck{Festes Gitter, erste Spalte immer Beschriftung. Nimmt rows, colsep
und grid unter denselben Namen wie die ZSFtable.}

\begin{valuegrid}[D-14 · valuegrid rows=tight][rows=tight, colsep=tight]{3}
  Mustertext & 12 & 24 & 36 \\
  Zweite Zeile & 20 & 40 & 60 \\
  Dritte Zeile & 30 & 60 & 90
\end{valuegrid}
\Zweck{Dieselben Masse-Regler wie an der Tabelle, hier zusammen gesetzt.}
